\section*{Capítulo \ref{ch:b2:neurons-synapses} --- Neurônios, potenciais de ação e sinapses}

\begin{solution}{exo:b2:neurons-synapses:1}
Os \omterm{def:b2:neurons-synapses:neuron}{dendritos} recebem as \omterm{def:b2:neurons-synapses:synapse}{sinapses}; o soma contém o núcleo e integra; o
\omterm{prop:b2:neurons-synapses:integration}{cone de implantação} decide; o \omterm{def:b2:neurons-synapses:neuron}{axônio} conduz; os terminais liberam o
transmissor. A \omterm{def:b2:neurons-synapses:neuron}{mielina} é a membrana enrolada de uma célula glial, quase
sem citoplasma; ela isola o \omterm{def:b2:neurons-synapses:neuron}{axônio}, de modo que o impulso salta entre os
nós e se propaga cem vezes mais depressa, a um custo menor.
\end{solution}

\begin{solution}{exo:b2:neurons-synapses:2}
Despolarização até o limiar; fase de subida --- os canais de sódio
dependentes de voltagem se abrem, retroalimentação positiva; pico perto
de $E_{\mathrm{Na}}$ --- os canais de sódio se inativam; fase de descida
--- os canais de potássio dependentes de voltagem se abrem;
hiperpolarização --- os canais de potássio ainda abertos, perto de
$E_{\mathrm K}$; recuperação à medida que eles se fecham. De tudo ou
nada porque a retroalimentação positiva ou vai até o fim ou não
acontece; unidirecional porque a membrana recém-percorrida está
refratária (canais de sódio inativados).
\end{solution}

\begin{solution}{exo:b2:neurons-synapses:3}
O impulso chega ao terminal; os canais de cálcio dependentes de voltagem
se abrem; o cálcio desencadeia a fusão das vesículas; o transmissor se
difunde através da fenda; liga-se aos receptores pós-sinápticos; os
canais se abrem (ou uma \omterm{prop:b2:cell-signalling:gpcr}{proteína G} é ativada); \omterm{def:b2:neurons-synapses:synapse}{potencial pós-sináptico};
o transmissor é removido por enzima ou transportador; a membrana da
vesícula é recuperada.
\end{solution}

\begin{solution}{exo:b2:neurons-synapses:4}
Elétrica: sem atraso, geralmente nos dois sentidos, sem amplificação,
sem inibição (a corrente apenas passa). Química: meio milissegundo,
unidirecional, amplificação (um impulso libera centenas de quanta),
inibição possível (canais de cloreto ou de potássio) e modificável.
\end{solution}

\begin{solution}{exo:b2:neurons-synapses:5}
A partir de $E = (61.5/z)\log_{10}(c_{\text{ext}}/c_{\text{int}})$:
potássio \qty{-89}{mV}; sódio \qty{+61}{mV}; cloreto, com $z =
-1$, $-61.5\log_{10} 11 = \qty{-64}{mV}$; cálcio, com $z = 2$,
$30.75\times 4.3 = \qty{+132}{mV}$. A \qty{-70}{mV}: o potássio sai, o
sódio entra, o cloreto sai ligeiramente (a membrana está abaixo de seu
equilíbrio), o cálcio entra com força.
\end{solution}

\begin{solution}{exo:b2:neurons-synapses:6}
$20 : 1$: $(20\times(-89) + 61)/21 = \qty{-82}{mV}$, o estado de
repouso, perto de $E_{\mathrm K}$. $1 : 20$: $(-89 + 20\times 61)/21 =
\qty{+54}{mV}$, o pico, perto de $E_{\mathrm{Na}}$. $50 : 1$: $(50\times
(-89) + 61)/51 = \qty{-86}{mV}$, a hiperpolarização, com canais de
potássio adicionais abertos.
\end{solution}

\begin{solution}{exo:b2:neurons-synapses:7}
$\sqrt{500} = 22$: \qty{22}{m/s}. Para chegar a \qty{100}{m/s}, um
\omterm{def:b2:neurons-synapses:neuron}{axônio} não mielinizado precisaria ter $d = 10^{4}$ \unit{\micro m}, um
centímetro. Um metro leva \qty{10}{ms} na fibra mielinizada,
\qty{45}{ms} no \omterm{def:b2:neurons-synapses:neuron}{axônio} de lula e um segundo inteiro numa fibra de
\qty{1}{\micro m}.
\end{solution}

\begin{solution}{exo:b2:neurons-synapses:8}
$m = \ln(200/74) = 1.0$; $P(1) = e^{-1} = 0.37$, $P(2) = 0.18$, $P(3)
= 0.06$. Resposta média $1.0\times 0.5 = \qty{0.5}{mV}$.
\end{solution}

\begin{solution}{exo:b2:neurons-synapses:9}
$1/\qty{2}{ms} = 500$ impulsos por segundo. A intensidade do toque é
codificada pela frequência dos impulsos (e pelo número de \omterm{def:b2:neurons-synapses:neuron}{neurônios}
recrutados); os próprios impulsos são idênticos.
\end{solution}

\begin{solution}{exo:b2:neurons-synapses:10}
$\lambda = \sqrt{R_m d/4R_i}$: \qty{0.5}{mm}, \qty{1.6}{mm},
\qty{11}{mm}. Com \omterm{def:b2:neurons-synapses:neuron}{mielina}, \qty{10}{\micro m}: $1.6\times\sqrt{200} =
\qty{22}{mm}$. Um internó de \qty{1}{mm} perde apenas $1 - e^{-1/22}
= \qty{4}{\%}$ do sinal, de modo que o nó seguinte é facilmente levado
ao limiar. Através de \qty{3}{mm} desmielinizados, o sinal cai para
$e^{-3/1.6} = \qty{15}{\%}$ --- de um impulso de \qty{100}{mV}, mais ou
menos o limiar de \qty{15}{mV}: a condução falha ou se torna incerta.
\end{solution}

\begin{solution}{exo:b2:neurons-synapses:11}
Tetrodotoxina: nenhuma corrente de sódio, nenhum impulso, nenhuma
transmissão. Tetraetilamônio: nenhuma corrente de potássio, impulso
prolongado, mais transmissor liberado por impulso. Sem cálcio: impulsos
normais, nenhuma liberação, nenhuma transmissão. Curare: liberação
normal, receptores bloqueados, nenhum potencial de placa motora ---
paralisia. Inibidor da colinesterase: a acetilcolina persiste, os
receptores permanecem abertos, o músculo se despolariza e depois não
consegue se repolarizar --- paralisia por excesso. Toxina botulínica:
as vesículas não conseguem se fundir, nenhuma liberação --- paralisia.
\end{solution}

\begin{solution}{exo:b2:neurons-synapses:12}
O atraso compra: a transmissão num único sentido; a amplificação (um
único impulso libera centenas de quanta e pode excitar uma célula
maior); a inibição, que uma junção elétrica não consegue fazer; e a
plasticidade, pois o número de quanta e de receptores pode mudar com o
uso, que é como os circuitos aprendem. As \omterm{def:b2:neurons-synapses:synapse}{sinapses} elétricas são usadas
onde a sincronia e a velocidade importam mais que o processamento: os
circuitos de fuga, o músculo cardíaco, algumas redes inibitórias.
\end{solution}

\begin{solution}{pb:b2:neurons-synapses:1}
\textbf{1.} $E_{\mathrm K} = 61.5\log_{10}(4/140) = \qty{-95}{mV}$;
$E_{\mathrm{Na}} = 61.5\log_{10}(145/12) = \qty{+67}{mV}$.
\textbf{2.} $(25\times(-95) + 67)/26 = \qty{-89}{mV}$.
\textbf{3.} $E_{\mathrm K} = 61.5\log_{10}(8/140) = \qty{-76}{mV}$;
$V = (25\times(-76) + 67)/26 = \qty{-71}{mV}$: mais perto do limiar,
portanto mais excitável de início --- e, se persistir, menos, pois os
canais de sódio se inativam no nível despolarizado.
\textbf{4.} Os gradientes se desfazem ao longo de horas, à medida que o
sódio entra e o potássio sai; o potencial deriva para zero; com a bomba
parada, os ânions impermeantes da célula atraem cloreto e água, e ela
incha.
\textbf{5.} Superfície por milímetro $\pi\times 15\times 10^{-6}\times
10^{-3} = \qty{4.7e-8}{m^2}$; $C = \qty{4.7e-10}{F}$; $Q = CV =
\qty{4.7e-11}{C}$: $2.9\times 10^{8}$ íons.
\textbf{6.} Volume por milímetro $\pi(7.5\times 10^{-6})^{2}\times
10^{-3} = \qty{1.8e-13}{m^3} = \qty{1.8e-10}{L}$, que contém $12\times
10^{-3}\times 1.8\times 10^{-10} = 2.1\times 10^{-12}$ mol, $1.3\times
10^{12}$ íons: uma mudança de $2\times 10^{-4}$.
\textbf{7.} $2.9\times 10^{8}/3 \approx 10^{8}$ ATP por milímetro por
impulso.
\textbf{8.} $0.9/90 = \qty{10}{ms}$ em cada sentido.
\textbf{9.} Nu: $\sqrt{1\times 15\times 10^{-6}/4} = \qty{1.9}{mm}$;
mielinizado: $\times\sqrt{250} = \qty{31}{mm}$.
\textbf{10.} Nu: $e^{-1.5/1.9} = 0.45$, \qty{45}{mV}; mielinizado:
$e^{-1.5/31} = 0.95$, \qty{95}{mV}. Ambos superam o limiar de
\qty{15}{mV}; mas o \omterm{def:b2:neurons-synapses:neuron}{axônio} nu precisa carregar toda a sua membrana no
caminho, o que o torna lento, ao passo que sob a \omterm{def:b2:neurons-synapses:neuron}{mielina} quase nada se
perde ou é carregado.
\textbf{11.} $\sqrt{15} = \qty{3.9}{m/s}$: \qty{0.23}{s} em cada
sentido, quase meio segundo para a ida e volta --- lento demais para um
reflexo que precisa corrigir um tropeço.
\textbf{12.} $e^{-4/1.9} = 0.12$: \qty{12}{mV} no nó distante, abaixo do
limiar --- o impulso é bloqueado e o reflexo se perde.
\textbf{13.} Atrás do impulso, os canais de sódio ficam inativados por
um ou dois milissegundos, de modo que a membrana recém-percorrida não
pode ser reexcitada e a onda só consegue avançar.
\textbf{14.} $m = \ln(300/111) = 1.0$.
\textbf{15.} $P(1) = 0.37$, $P(2) = 0.18$, $P(3) = 0.06$: cerca de 110,
55 e 18 das 300 tentativas.
\textbf{16.} $1.0\times 0.4 = \qty{0.4}{mV}$.
\textbf{17.} $P(0) = e^{-60} \approx 10^{-26}$: nunca; resposta média
de \qty{24}{mV}, acima do limiar de \qty{15}{mV}: uma única \omterm{def:b2:neurons-synapses:synapse}{sinapse}
desse tipo faz o \omterm{def:b2:neurons-synapses:neuron}{neurônio} motor disparar.
\textbf{18.} Vinte \omterm{def:b2:neurons-synapses:synapse}{sinapses} ativas ao mesmo tempo somam suas correntes
(somação espacial), de modo sublinear à medida que o potencial se
aproxima do potencial de reversão dos canais; mesmo com $m = 1$ cada
uma, dariam \qty{8}{mV}, e com o $m$ normal muito mais que o limiar.
\textbf{19.} Abrir canais de cloreto com $E_{\mathrm{Cl}} = V_{\text{repouso}}$
acrescenta condutância sem mover o potencial; a corrente excitatória
passa então por uma membrana com mais fugas e produz uma despolarização
menor --- é a inibição por derivação.
\textbf{20.} $10 + 0.6 + 10 + 0.6 + 5 = \qty{26}{ms}$, contra
\qty{30}{ms} medidos.
\textbf{21.} A ativação do próprio receptor no fuso muscular, a
propagação do \omterm{prop:b2:neurons-synapses:ap}{potencial de ação} da \omterm{def:b2:cell-differentiation:myogenesis}{fibra muscular} ao longo da fibra
(metros por segundo ao longo de centímetros) e a latência mecânica antes
que a força apareça.
\textbf{22.} Cada \omterm{def:b2:neurons-synapses:synapse}{sinapse} acrescenta meio milissegundo e um interneurônio
acrescenta uma célula; mas os interneurônios permitem inverter o sinal
(inibindo o antagonista), combiná-lo com outros e deixá-lo ser modulado
pelo encéfalo --- um reflexo que pode ser moldado.
\textbf{23.} $90 - 6.2 = \qty{84}{ms}$ de condução para \qty{1.8}{m}:
cerca de \qty{21}{m/s}.
\textbf{24.} Metade da corrente de entrada em cada voltagem: o limiar
sobe, o impulso fica menor e mais lento, a margem de segurança em cada
nó diminui e o impulso pode falhar; o reflexo enfraquece ou desaparece.
\textbf{25.} \omterm{thm:b2:neurons-synapses:chord}{Potencial de repouso} de \qty{-89}{mV}; \qty{10}{ms} em cada
sentido; $m = 1.0$; latência calculada de \qty{26}{ms}.
\end{solution}
